\documentclass{note}

\title{Expression problem}
\date{2026-04-26}
\modified{2026-04-26}

\begin{document}

The \emph{expression problem} observes the inherent conflict between adding new types and new operations to open systems.

Let's use the classical \code{Shape} types as an example.
We start with two types (\code{Rectangle} and \code{Circle}) and two operations (\code{area} and \code{perimeter}).

\section{adding-types}{Adding types}

Object-oriented paradigm facilitates adding new types (e.g., \code{Triangle}).

\begin{code}[ruby]
class Rectangle
  def area
    @width * @height
  end

  def perimeter
    2 * (@width + @height)
  end
end

class Circle
  def area
    Math::PI * @radius * @radius
  end

  def perimeter
    2 * Math::PI * @radius
  end
end
\end{code}

Adding a new method (e.g., \code{vertex\_count}) requires changes in every class.

\section{adding-functions}{Adding functions}

Functional programming with algebraic data types makes adding new operations easy.

\begin{code}[ocaml]
type shape =
  | Rectangle of float * float
  | Circle of float

let area = function
  | Rectangle (w, h) -> w *. h
  | Circle r -> Float.pi *. r *. r

let perimeter = function
  | Rectangle (w, h) -> 2.0 *. (w +. h)
  | Circle r -> 2.0 *. Float.pi *. r
\end{code}

Adding new shapes requires changes in all existing functions.

The visitor design pattern in \textsc{oop} trades one type of extensibility for another,
mimicking the functional style and inheriting its limitations.

\begin{code}[ruby]
class Visitor
  def visit\_rectangle(w, h)
    raise NotImplementedError
  end

  def visit\_circle(r)
    raise NotImplementedError
  end
end

class Rectangle
  def visit(v)
    return v.visit\_rectangle(@w, @h)
  end
end

class Circle
  def visit(v)
    return v.visit\_circle(@r)
  end
end

class Area
  def visit\_rectangle(w, h)
    return w * h
  end

  def visit\_circle(r)
    return Math::PI * r * r
  end
end
\end{code}
\end{document}
